\section{A General Equilibrium}

\begin{frame}{From Partial to General Equilibrium}
    \mytext{
    Microeconomics demonstrated the equilibrium of supply and demand for individual markets: partial equlibria.

    We have also seen that goods can be complementary or substitutes. The price of one good (equilibrium on that market's good) influences the price of another (cross-price elasticities).

    So the questions are: how does one market influence another? What can we say about an equilibrium on multiple markets, a general equilibrium?
    }
\end{frame}

\begin{frame}{Say's Law}
    \mytext{
    Every demand equals supply, as we produce to have something to exchange for what we desire.

    Every supply equals demand, as everything produced will at some point be exchanged (or self-consumed).

    Production is the source of demand.
    }
\end{frame}

\begin{frame}{Walras' Law}
    \mytext{
        If one market is in disequilibrium, another market must also be in disequilibrium.

        The sum of all derivations from equilibrium is zero.

        It follows that if $n-1$ markets are in equilibrium, market $n$ is also in equilibrium.
    }
\end{frame}

\begin{frame}{Patinkin's Model of Money, Interest, and Prices}
    \showmytikz[-1,1]{../shared/tikz/tikz-threemarkets}
    \mycomment{
    Think of two (partly substitutable) goods $X,Y$ and a denominator $Z$, like fish, coconut and shells. 

    As the price of one good increases, the price of the other increases as well (positive cross-price effect, substitutes). 

    There is a price ratio for $X,Y$  where supply and demand for $X$ are in equilibrium.

    \bigskip 
    
        All this follows Krugman: \href{https://web.mit.edu/krugman/www/islm.html}{https://web.mit.edu/krugman/www/islm.html}
    
    }
\end{frame}

\begin{frame}{Patinkin's Model of Money, Interest, and Prices}
    \showmytikz[-2,1]{../shared/tikz/tikz-threemarkets}
    \mycomment{
    The same applies for $Y$. 
    
    Both slope upwards since they are substitutes; their cross-price elasticity is positive.
    
    Now in the case of only two goods, both markets must be in equilibrium simultaneously. Thus, there will be a certain number  $X^*,Y^*$ sold and bought for a specific price.
    }
\end{frame}

\begin{frame}{Patinkin's Model of Money, Interest, and Prices}
    \showmytikz[1]{../shared/tikz/tikz-threemarkets}
    \mycomment{

    
    Walras' law states: when $n-1$ markets are in balance, the remaining market must also be in balance.

    Thus, we can also have a third good $Z$, in our example the shells. 
    
    % Think of two (partly compementary) goods $X,Y$ and a denominator $Z$, like fish, coconut and shells.
    }
\end{frame}

\begin{frame}{Patinkin's Model of Money, Interest, and Prices}
    \showmytikz[2]{../shared/tikz/tikz-threemarkets}
    \mycomment{
    We can think of these three goods as goods `in general' ($X$), bonds `in general' ($Y$) and money ($Z$). (Or consumption vs. savings.)

    So money can be either held, or spent on goods (consumed), or invested in bonds.

    The price of bonds is typically expressed by the interest rate; so the price axis `price of bonds' is now inverse. Let us flip it so that the axis shows the interest rate.
    }
\end{frame}

\begin{frame}{Patinkin's Model of Money, Interest, and Prices}
    \showmytikz[3]{../shared/tikz/tikz-threemarkets}
    \mycomment{
    Now we have a general equilibrium model of three markets: goods, money, and bonds.
    }
\end{frame}

\begin{frame}{Patinkin's Model of Money, Interest, and Prices}
    \showmytikz[3,4]{../shared/tikz/tikz-threemarkets}
    \mycomment{
        If the market for goods remains stable, but there is more supply of money, interest rates will go down and the price of goods will increase.

        This is equal to a monetary stimulus: the government (central bank) buys bonds and `sells' newly printed money. Cheper money reduces the interest rate at higher output.
    }
\end{frame}

\begin{frame}{Patinkin's Model of Money, Interest, and Prices}
    \showmytikz[3,5]{../shared/tikz/tikz-threemarkets}
    \mycomment{
        If money supply remains stable, an increase of interest rates will result in a higher price for goods.

        Or say it this way: an investment boom (bonds curve shift) leads to higher output (goods curve shift) at higher interest rates ($i_2>i^*$) and higher inflation ($P_2>P^*$).

        This will also happen when the government sells bonds and buys goods (fiscal stimulus).
    }
\end{frame}

% \begin{frame}{Patinkin's Model of Money, Interest, and Prices}
%     \showmytikz[3,6]{../shared/tikz/tikz-threemarkets}
%     \mycomment{
%         A stable supply of bonds means more money will push the price of goods down at lower interest rates. 
        
        
%     }
% \end{frame}